\chapter{Двухкопийная кратность спинорного накрытия: окончательный критерий остановки}

Предыдущая проверка оставила одну лазейку: возможно, требуемая внутренняя
двойка массы $n\cdot\sigma$ уже скрыта в тензорном квадрате, SWAP или
внешней степени носителях Тома VI. Ретроспективный атлас уточнил численную
мишень через
\begin{equation}
 \Delta_{\rm crit}^2=\frac14=\frac6{24},
 \qquad
 R_{\rm crit}=\frac1{24}\operatorname{diag}(16,4,4).
\end{equation}
Настоящая проверка выясняет не совпадение размерностей, а существование
канонического пространства кратности $\mathbb C^2$, на котором полный
матричный алгебраический фактор $M_2(\mathbb C)$ коммутирует с
физической калибровочной алгеброй и статистическим проектором.

\section{Семейный тензорный квадрат не содержит кратности два}

Диагональное действие вращений на уже существующем семейном носителе
имеет стандартное разложение
\begin{equation}
 \mathbb C^3\otimes\mathbb C^3
 \simeq V_0\oplus V_1\oplus V_2,
 \qquad
 \dim(V_0,V_1,V_2)=(1,3,5).
 \label{eq:v6-two-copy-so3-decomposition}
\end{equation}
Машинная диагонализация полного казимира дала собственные значения
$0,2,6$ с кратностями $1,3,5$. Каждое неприводимое представление входит
ровно один раз. Поэтому коммутант диагонального $SO(3)$ имеет размерность
\begin{equation}
 1^2+1^2+1^2=3
\end{equation}
и является абелевой алгеброй спектральных проекторов на три слагаемых.
В нём нет простого блока $M_2(\mathbb C)$.

Это является конкретной формой двойного коммутантного принципа
Шура--Вейля. Современная higher-spin формулировка дана в
\path{arXiv:2008.06038}; она также подчёркивает, что физическая матричная
алгебра возникает именно на пространствах кратности, а не из одной
размерности неприводимого сектора.

\begin{proposition}[Отсутствие семейного Pauli-фактора]
Тензорный квадрат векторного семейного триплета содержит спины ноль, один
и два без повторений. Следовательно, он не порождает каноническую
комплексную двойку, независимую от $SO(3)$-действия.
\end{proposition}

\section{Два порядка обмена исчезают после статистического выбора}

До факторизации можно записать формальную пару порядков
\begin{equation}
 \operatorname{span}\{|LR\rangle,|RL\rangle\}\simeq\mathbb C^2.
\end{equation}
Но перестановка действует матрицей $\sigma_x$. Её коммутант равен
\begin{equation}
 \operatorname{span}\{I,\sigma_x\},
\end{equation}
то есть двумерной абелевой алгебре. Нормы коммутаторов перестановки с
$(\sigma_x,\sigma_y,\sigma_z)$ равны
\begin{equation}
 (0,2\sqrt2,2\sqrt2).
\end{equation}
После выбора симметричного или антисимметричного сектора остаётся одна
комплексная координата, а не дублет. Недиагональные Pauli-генераторы
меняют статистический сектор и потому не являются физическим внутренним
$SU(2)$.

\section{Тензорные произведения физического пакета}

Проверено точное представительное содержание
\begin{equation}
 H_{15}=Q_L\oplus L_L\oplus u_R\oplus d_R\oplus e_R
\end{equation}
с принятыми квантовыми числами
\begin{equation}
 (3,2,1/6),(1,2,-1/2),(3,1,2/3),(3,1,-1/3),(1,1,-1).
\end{equation}
В прямом квадрате целевые кратности этих пяти блоков равны
\begin{equation}
 (2,0,0,4,1),
\end{equation}
то есть отсутствуют $L_L$ и $u_R$. В смешанном квадрате
$H_{15}\otimes H_{15}^*$ получено
\begin{equation}
 (1,3,2,0,1),
\end{equation}
то есть отсутствует $d_R$. Сопряжённый порядок даёт тот же дефицит.

Следовательно, тензорные произведения содержат отдельные знакомые
представления, как и обычные составные или юкавские каналы, но не
содержат равномерной равнозарядной копии всего $H_{15}$. Тем более они не
дают две такие копии, на которых Pauli-действие коммутировало бы со всей
калибровочной алгеброй.

\begin{nogocriterion}[Составной канал не равен внутренней кратности]
Наличие отдельных блоков Стандартной модели в $H_{15}^{\otimes2}$ не
позволяет читать их как второй экземпляр фермионного пакета. Необходима
одна общая интертвинерная копия каждого из пяти блоков с одинаковой
кратностью и согласованной градуировкой.
\end{nogocriterion}

\section{Почему ранговый мост атласа не спасает дублет}

В разложении
\begin{equation}
 24=8_C+3_W+1_Y+6_X+6_{\bar X}
\end{equation}
ранг четыре означает $W\oplus Y$. Но ранг шестнадцать есть дополнение
цветового блока и уже содержит тот же $W\oplus Y$. Для соответствующих
проекторов
\begin{equation}
 P_4P_{16}=P_4.
\end{equation}
Если два написанных ранга четыре адреса считать двумя кандидатами копии, их
матрица перекрытий равна
\begin{equation}
 \begin{pmatrix}4&4\\4&4\end{pmatrix},
\end{equation}
имеет собственные значения $(0,8)$ и ранг один. Следовательно, запись
$(16,4,4)/24$ повторяет один проектор численно, но не создаёт два
ортогональных подпространства.

Ранги $X$ и $\bar X$ оба равны шести, однако это сопряжённые калибровочные блоки.
Для зарядового прокси $\sigma_z\otimes I_6$ нормы коммутаторов с Pauli-
тройкой равны $(4\sqrt3,4\sqrt3,0)$. Поэтому смешивать их можно только
диагональным $U(1)$, как и прежнюю KO6-пару.

Тождество $\Delta^2=6/24$ остаётся точной спектральной подсказкой, но
равенство числа рангу не выводит внутренний носитель спинорного накрытия.

\section{Вердикт}

Последняя лазейка текущего конечного родителя закрыта:
\begin{equation}
 \boxed{
 \begin{gathered}
  \text{носители тензорного квадрата, оператора SWAP и внешних степеней}\\
  \text{не дают физический }\mathbb C^2_{\rm twist}.
 \end{gathered}}
\end{equation}
Это не уничтожает выведенное бозонное поле $Q$, конечный составной дефект,
хопфову линию или граничный класс $\pm15$. Закрывается более сильное
утверждение, что пятнадцать локализованных фермионных мод уже следуют из
$H_{15}/M_{35}$ без нового представительного ресурса.

Переоткрытие возможно только честным изменением архитектуры:
\begin{enumerate}
 \item добавить равнозарядный модуль кратности с новым следом,
 аномальным и first-order аудитом;
 \item либо объявить структуру типа
 $\operatorname{Spin}^h=(\operatorname{Spin}\times Sp(1))/\mathbb Z_2$,
 связывающую пространственный spin с независимым кватернионным twist;
 такой маршрут является новой версией, а не скрытым свойством старой.
\end{enumerate}
Литературный язык второго варианта дан в
\path{arXiv:2008.04934}. Теорема Каллиаса по-прежнему вычисляет индекс
уже заданного Дирака--Хиггса оператора и не создаёт отсутствующую
представительную кратность; см. Anghel,
\path{Commun.Math.Phys.128(1990)77-97}.

Следующий маршрут Тома VI должен вернуться к безусловно выведенному
бозонному дефекту и спросить, какие наблюдаемые поля и материальные
квантовые числа следуют из него без условной фермионной интерпретации.

\begin{protocol}
\begin{enumerate}
 \item $3\otimes3=1\oplus3\oplus5$: \textbf{точно};
 \item кратности неприводимых семейных секторов: \textbf{все единичны};
 \item $M_2$ в семейном коммутанте: \textbf{отсутствует};
 \item формальная пара $LR/RL$: \textbf{есть до факторизации};
 \item после статистического проектора: \textbf{одномерные сектора};
 \item полный $H_{15}$ в $H_{15}\otimes H_{15}$: \textbf{нет};
 \item полный $H_{15}$ в смешанном квадрате: \textbf{нет};
 \item повтор ранга четыре адреса атласа: \textbf{один и тот же проектор};
 \item $X/\bar X$ как равнозарядный дублет: \textbf{нет};
 \item физический носитель спинорного накрытия из текущего родителя:
 \textbf{не выведен};
 \item бозонный дефект и класс $\pm15$: \textbf{сохраняются};
 \item фермионная ветвь спинорного накрытия текущей архитектуры: \textbf{закрыта};
 \item следующая проверка: \textbf{идентификация бозонного дефектного поля}.
\end{enumerate}
\end{protocol}

Машинный сертификат сохранён в
\path{s2t_v6_two_copy_spin_cover_multiplicity_gate_results.json}.